\documentclass[10pt,a4paper]{amsart}
\usepackage{amsmath,amssymb,amsthm,graphicx}
% Retained preamble for the Stacks Project chapter snapshot in this directory.
% The chapter's own preamble.tex is a copy of the shared Stacks Project preamble; the block
% below is sliced MECHANICALLY, line for line, from it. Nothing was added, removed or
% redefined; the wrapper class, the xr-hyper/hyperref wiring and the externaldocument list
% are omitted because the assembler supplies the document class.
%
% lines taken from preamble.tex: 3, 5, 9, 11, 12, 156, 157, 158, ... (51 lines total)
%
% package coverage restored from the chapter's own preamble (2026-10-13): xy, xy-2cell

\usepackage[all]{xy}\usepackage{verbatim}
\xyoption{2cell}
\UseAllTwocells

\newenvironment{reference}{\comment}{\endcomment}
\newenvironment{slogan}{\comment}{\endcomment}
\newenvironment{history}{\comment}{\endcomment}
\theoremstyle{plain}
\newtheorem{theorem}[subsection]{Theorem}
\newtheorem{proposition}[subsection]{Proposition}
\newtheorem{lemma}[subsection]{Lemma}
\theoremstyle{definition}
\newtheorem{definition}[subsection]{Definition}
\newtheorem{example}[subsection]{Example}
\newtheorem{exercise}[subsection]{Exercise}
\newtheorem{situation}[subsection]{Situation}
\theoremstyle{remark}
\newtheorem{remark}[subsection]{Remark}
\newtheorem{remarks}[subsection]{Remarks}
\numberwithin{equation}{subsection}
\def\lim{\mathop{\mathrm{lim}}\nolimits}
\def\colim{\mathop{\mathrm{colim}}\nolimits}
\def\Spec{\mathop{\mathrm{Spec}}}
\def\Hom{\mathop{\mathrm{Hom}}\nolimits}
\def\Ext{\mathop{\mathrm{Ext}}\nolimits}
\def\SheafHom{\mathop{\mathcal{H}\!\mathit{om}}\nolimits}
\def\SheafExt{\mathop{\mathcal{E}\!\mathit{xt}}\nolimits}
\def\Sch{\mathit{Sch}}
\def\Mor{\mathop{\mathrm{Mor}}\nolimits}
\def\Ob{\mathop{\mathrm{Ob}}\nolimits}
\def\Sh{\mathop{\mathit{Sh}}\nolimits}
\def\NL{\mathop{N\!L}\nolimits}
\def\CH{\mathop{\mathrm{CH}}\nolimits}
\def\proetale{{pro\text{-}\acute{e}tale}}
\def\etale{{\acute{e}tale}}
\def\QCoh{\mathit{QCoh}}
\def\Ker{\mathop{\mathrm{Ker}}}
\def\Im{\mathop{\mathrm{Im}}}
\def\Coker{\mathop{\mathrm{Coker}}}
\def\Coim{\mathop{\mathrm{Coim}}}
\DeclareMathSymbol{\boxtimes}{\mathbin}{AMSa}{"02}
\def\QCohstack{\mathcal{QC}\!\mathit{oh}}
\def\Cohstack{\mathcal{C}\!\mathit{oh}}
\def\Spacesstack{\mathcal{S}\!\mathit{paces}}
\def\Quotfunctor{\mathrm{Quot}}
\def\Hilbfunctor{\mathrm{Hilb}}
\def\Curvesstack{\mathcal{C}\!\mathit{urves}}
\def\Polarizedstack{\mathcal{P}\!\mathit{olarized}}
\def\Complexesstack{\mathcal{C}\!\mathit{omplexes}}
\def\Pic{\mathop{\mathrm{Pic}}\nolimits}
\def\Picardstack{\mathcal{P}\!\mathit{ic}}
\def\Picardfunctor{\mathrm{Pic}}
\def\Deformationcategory{\mathcal{D}\!\mathit{ef}}
\usepackage{cleveref}
\crefname{originallemma}{original lemma}{original lemmas}
\Crefname{originallemma}{Original lemma}{Original lemmas}
\makeatletter
\newlabel{algebraic-lemma-map-space-into-stack}{{04T4}{source}{}{sourcepointer-1}{}}
\newlabel{algebraic-lemma-map-space-into-stack@cref}{{[originallemma][0][]04T4}{[1][0][]source}{}{}{}}
\makeatother
\crefname{originallemma}{original lemma}{original lemmas}
\Crefname{originallemma}{Original lemma}{Original lemmas}
\makeatletter
\newlabel{lemma-etale-local-lifts-isomorphic}{{lemma-etale-local-lifts-isomorphic}{source}{}{sourcepointer-2}{}}
\newlabel{lemma-etale-local-lifts-isomorphic@cref}{{[originallemma][0][]lemma-etale-local-lifts-isomorphic}{[1][0][]source}{}{}{}}
\makeatother
\crefname{originallemma}{original lemma}{original lemmas}
\Crefname{originallemma}{Original lemma}{Original lemmas}
\makeatletter
\newlabel{lemma-gerbe-of-lifts-fibred}{{lemma-gerbe-of-lifts-fibred}{source}{}{sourcepointer-3}{}}
\newlabel{lemma-gerbe-of-lifts-fibred@cref}{{[originallemma][0][]lemma-gerbe-of-lifts-fibred}{[1][0][]source}{}{}{}}
\makeatother
\crefname{originallemma}{original lemma}{original lemmas}
\Crefname{originallemma}{Original lemma}{Original lemmas}
\makeatletter
\newlabel{remark-gerbe-of-lifts}{{remark-gerbe-of-lifts}{source}{}{sourcepointer-4}{}}
\newlabel{remark-gerbe-of-lifts@cref}{{[originallemma][0][]remark-gerbe-of-lifts}{[1][0][]source}{}{}{}}
\makeatother
\crefname{originallemma}{original lemma}{original lemmas}
\Crefname{originallemma}{Original lemma}{Original lemmas}
\makeatletter
\newlabel{sites-section-glueing-sheaves}{{04TP}{source}{}{sourcepointer-5}{}}
\newlabel{sites-section-glueing-sheaves@cref}{{[originallemma][0][]04TP}{[1][0][]source}{}{}{}}
\makeatother
\crefname{originallemma}{original lemma}{original lemmas}
\Crefname{originallemma}{Original lemma}{Original lemmas}
\makeatletter
\newlabel{spaces-more-groupoids-lemma-composition-is-addition}{{0CKF}{source}{}{sourcepointer-6}{}}
\newlabel{spaces-more-groupoids-lemma-composition-is-addition@cref}{{[originallemma][0][]0CKF}{[1][0][]source}{}{}{}}
\makeatother
\crefname{originallemma}{original lemma}{original lemmas}
\Crefname{originallemma}{Original lemma}{Original lemmas}
\makeatletter
\newlabel{spaces-more-groupoids-section-groupoid-sections}{{0CKB}{source}{}{sourcepointer-7}{}}
\newlabel{spaces-more-groupoids-section-groupoid-sections@cref}{{[originallemma][0][]0CKB}{[1][0][]source}{}{}{}}
\makeatother
\crefname{originallemma}{original lemma}{original lemmas}
\Crefname{originallemma}{Original lemma}{Original lemmas}
\makeatletter
\newlabel{spaces-more-morphisms-definition-conormal-sheaf}{{04CN}{source}{}{sourcepointer-8}{}}
\newlabel{spaces-more-morphisms-definition-conormal-sheaf@cref}{{[originallemma][0][]04CN}{[1][0][]source}{}{}{}}
\makeatother
\crefname{originallemma}{original lemma}{original lemmas}
\Crefname{originallemma}{Original lemma}{Original lemmas}
\makeatletter
\newlabel{spaces-more-morphisms-lemma-action-by-derivations-etale-localization}{{0CK8}{source}{}{sourcepointer-9}{}}
\newlabel{spaces-more-morphisms-lemma-action-by-derivations-etale-localization@cref}{{[originallemma][0][]0CK8}{[1][0][]source}{}{}{}}
\makeatother
\crefname{originallemma}{original lemma}{original lemmas}
\Crefname{originallemma}{Original lemma}{Original lemmas}
\makeatletter
\newlabel{spaces-more-morphisms-lemma-action-sheaf}{{061C}{source}{}{sourcepointer-10}{}}
\newlabel{spaces-more-morphisms-lemma-action-sheaf@cref}{{[originallemma][0][]061C}{[1][0][]source}{}{}{}}
\makeatother
\crefname{originallemma}{original lemma}{original lemmas}
\Crefname{originallemma}{Original lemma}{Original lemmas}
\makeatletter
\newlabel{spaces-more-morphisms-lemma-smooth-omega-finite-locally-free}{{0CK5}{source}{}{sourcepointer-11}{}}
\newlabel{spaces-more-morphisms-lemma-smooth-omega-finite-locally-free@cref}{{[originallemma][0][]0CK5}{[1][0][]source}{}{}{}}
\makeatother
\crefname{originallemma}{original lemma}{original lemmas}
\Crefname{originallemma}{Original lemma}{Original lemmas}
\makeatletter
\newlabel{spaces-more-morphisms-remark-another-special-case}{{0CK7}{source}{}{sourcepointer-12}{}}
\newlabel{spaces-more-morphisms-remark-another-special-case@cref}{{[originallemma][0][]0CK7}{[1][0][]source}{}{}{}}
\makeatother
\crefname{originallemma}{original lemma}{original lemmas}
\Crefname{originallemma}{Original lemma}{Original lemmas}
\makeatletter
\newlabel{spaces-more-morphisms-theorem-topological-invariance}{{05ZH}{source}{}{sourcepointer-13}{}}
\newlabel{spaces-more-morphisms-theorem-topological-invariance@cref}{{[originallemma][0][]05ZH}{[1][0][]source}{}{}{}}
\makeatother
\crefname{originallemma}{original lemma}{original lemmas}
\Crefname{originallemma}{Original lemma}{Original lemmas}
\makeatletter
\newlabel{spaces-properties-lemma-characterize-quasi-coherent}{{03M0}{source}{}{sourcepointer-14}{}}
\newlabel{spaces-properties-lemma-characterize-quasi-coherent@cref}{{[originallemma][0][]03M0}{[1][0][]source}{}{}{}}
\makeatother
\crefname{originallemma}{original lemma}{original lemmas}
\Crefname{originallemma}{Original lemma}{Original lemmas}
\makeatletter
\newlabel{spaces-properties-lemma-properties-quasi-coherent}{{03M1}{source}{}{sourcepointer-15}{}}
\newlabel{spaces-properties-lemma-properties-quasi-coherent@cref}{{[originallemma][0][]03M1}{[1][0][]source}{}{}{}}
\makeatother
\crefname{originallemma}{original lemma}{original lemmas}
\Crefname{originallemma}{Original lemma}{Original lemmas}
\makeatletter
\newlabel{stacks-lemma-gerbe-abelian-auts}{{0CJY}{source}{}{sourcepointer-16}{}}
\newlabel{stacks-lemma-gerbe-abelian-auts@cref}{{[originallemma][0][]0CJY}{[1][0][]source}{}{}{}}
\makeatother
\title{Automorphism sheaves of the first-order lifting gerbe are quasi-coherent}
\author{The Stacks Project Authors}
\date{stacks-project:stacks-more-morphisms.tex:local-snapshot-20261012; GNU FDL 1.2 or later}
\begin{document}
\maketitle
\noindent\textit{Editorial note.} Statement span is lines 1322-1334 of sources/stacks-stacks-more-morphisms/pinned-original-source.tex; proof span is lines 1336-1495, the complete author proof (158 body lines).\par\medskip
\noindent\textit{Editorial note.} Navigation stubs: 16 cited-but-not-included labels are supplied as non-load-bearing pointers with their pinned tags where the snapshot tag index has them.\par\medskip
\noindent\textit{Editorial note.} extra\_packages is empty on purpose: the retained preamble already loads xy with options and a second declaration clashes in the assembler wrapper.\par\medskip
\noindent\textit{Editorial note.} Text is retained verbatim, including the author's own typos and placeholders.\par\medskip
\noindent\textit{Editorial note.} Queue index 3 of queue\_stacks-stacks-more-morphisms.json (pb=6696, stmt 1322-1334, proof 1336-1495).\par\medskip
\noindent\textit{Editorial note.} stacks-stacks-more-morphisms wave smm01 (2026-10-14): fresh adjudication accept per lane-s3/smm\_judgments.md; pool per Chief 10:5x open-pool ruling (net 59 with-proof rows, bloodline lane-s3 by anchor audit, v2 true-key namespace stacks:stacks-more-morphisms.tex, row lemma-deform-property-fp-over-ft avoided on cross-chapter key collision); 60 with-proof rows judged 45 accept / 14 reject / 1 avoid; patch-sd-envs contract spans scanned.\par\medskip
\noindent\textit{Allowed standard prerequisites (editorial).}\begin{itemize}
\item Explicit prerequisite boundary. The cited inputs are prerequisites; the argument in the retained span is not replaced by any cited result.
\end{itemize}\medskip
\ifdefined\hypertarget\hypertarget{sourcepointer-1}{}\fi
\noindent\textit{Source reference.} Original lemma 04T4: sources/stacks-stacks-more-morphisms/pinned-original-source.tex, source lines 1495--1495. Cited by the retained statement or proof but not part of the collected spans; supplied only as a navigation pointer. The argument does not depend on its content.\par\medskip
\ifdefined\hypertarget\hypertarget{sourcepointer-2}{}\fi
\noindent\textit{Source reference.} Original lemma lemma-etale-local-lifts-isomorphic: sources/stacks-stacks-more-morphisms/pinned-original-source.tex, source lines 1211--1226. Cited by the retained statement or proof but not part of the collected spans; supplied only as a navigation pointer. The argument does not depend on its content.\par\medskip
\ifdefined\hypertarget\hypertarget{sourcepointer-3}{}\fi
\noindent\textit{Source reference.} Original lemma lemma-gerbe-of-lifts-fibred: sources/stacks-stacks-more-morphisms/pinned-original-source.tex, source lines 964--967. Cited by the retained statement or proof but not part of the collected spans; supplied only as a navigation pointer. The argument does not depend on its content.\par\medskip
\ifdefined\hypertarget\hypertarget{sourcepointer-4}{}\fi
\noindent\textit{Source reference.} Original lemma remark-gerbe-of-lifts: sources/stacks-stacks-more-morphisms/pinned-original-source.tex, source lines 875--903. Cited by the retained statement or proof but not part of the collected spans; supplied only as a navigation pointer. The argument does not depend on its content.\par\medskip
\ifdefined\hypertarget\hypertarget{sourcepointer-5}{}\fi
\noindent\textit{Source reference.} Original lemma 04TP: sources/stacks-stacks-more-morphisms/pinned-original-source.tex, source lines 1495--1495. Cited by the retained statement or proof but not part of the collected spans; supplied only as a navigation pointer. The argument does not depend on its content.\par\medskip
\ifdefined\hypertarget\hypertarget{sourcepointer-6}{}\fi
\noindent\textit{Source reference.} Original lemma 0CKF: sources/stacks-stacks-more-morphisms/pinned-original-source.tex, source lines 1495--1495. Cited by the retained statement or proof but not part of the collected spans; supplied only as a navigation pointer. The argument does not depend on its content.\par\medskip
\ifdefined\hypertarget\hypertarget{sourcepointer-7}{}\fi
\noindent\textit{Source reference.} Original lemma 0CKB: sources/stacks-stacks-more-morphisms/pinned-original-source.tex, source lines 1495--1495. Cited by the retained statement or proof but not part of the collected spans; supplied only as a navigation pointer. The argument does not depend on its content.\par\medskip
\ifdefined\hypertarget\hypertarget{sourcepointer-8}{}\fi
\noindent\textit{Source reference.} Original lemma 04CN: sources/stacks-stacks-more-morphisms/pinned-original-source.tex, source lines 1495--1495. Cited by the retained statement or proof but not part of the collected spans; supplied only as a navigation pointer. The argument does not depend on its content.\par\medskip
\ifdefined\hypertarget\hypertarget{sourcepointer-9}{}\fi
\noindent\textit{Source reference.} Original lemma 0CK8: sources/stacks-stacks-more-morphisms/pinned-original-source.tex, source lines 1495--1495. Cited by the retained statement or proof but not part of the collected spans; supplied only as a navigation pointer. The argument does not depend on its content.\par\medskip
\ifdefined\hypertarget\hypertarget{sourcepointer-10}{}\fi
\noindent\textit{Source reference.} Original lemma 061C: sources/stacks-stacks-more-morphisms/pinned-original-source.tex, source lines 1495--1495. Cited by the retained statement or proof but not part of the collected spans; supplied only as a navigation pointer. The argument does not depend on its content.\par\medskip
\ifdefined\hypertarget\hypertarget{sourcepointer-11}{}\fi
\noindent\textit{Source reference.} Original lemma 0CK5: sources/stacks-stacks-more-morphisms/pinned-original-source.tex, source lines 1495--1495. Cited by the retained statement or proof but not part of the collected spans; supplied only as a navigation pointer. The argument does not depend on its content.\par\medskip
\ifdefined\hypertarget\hypertarget{sourcepointer-12}{}\fi
\noindent\textit{Source reference.} Original lemma 0CK7: sources/stacks-stacks-more-morphisms/pinned-original-source.tex, source lines 1495--1495. Cited by the retained statement or proof but not part of the collected spans; supplied only as a navigation pointer. The argument does not depend on its content.\par\medskip
\ifdefined\hypertarget\hypertarget{sourcepointer-13}{}\fi
\noindent\textit{Source reference.} Original lemma 05ZH: sources/stacks-stacks-more-morphisms/pinned-original-source.tex, source lines 1495--1495. Cited by the retained statement or proof but not part of the collected spans; supplied only as a navigation pointer. The argument does not depend on its content.\par\medskip
\ifdefined\hypertarget\hypertarget{sourcepointer-14}{}\fi
\noindent\textit{Source reference.} Original lemma 03M0: sources/stacks-stacks-more-morphisms/pinned-original-source.tex, source lines 1495--1495. Cited by the retained statement or proof but not part of the collected spans; supplied only as a navigation pointer. The argument does not depend on its content.\par\medskip
\ifdefined\hypertarget\hypertarget{sourcepointer-15}{}\fi
\noindent\textit{Source reference.} Original lemma 03M1: sources/stacks-stacks-more-morphisms/pinned-original-source.tex, source lines 1495--1495. Cited by the retained statement or proof but not part of the collected spans; supplied only as a navigation pointer. The argument does not depend on its content.\par\medskip
\ifdefined\hypertarget\hypertarget{sourcepointer-16}{}\fi
\noindent\textit{Source reference.} Original lemma 0CJY: sources/stacks-stacks-more-morphisms/pinned-original-source.tex, source lines 1495--1495. Cited by the retained statement or proof but not part of the collected spans; supplied only as a navigation pointer. The argument does not depend on its content.\par\medskip
\begin{lemma}
\label{lemma-gerbe-of-lifts-first-order}
In Remark \ref{remark-gerbe-of-lifts} assume $\mathcal{X} \subset \mathcal{X}'$
is a first order thickening. Then
\begin{enumerate}
\item the automorphism sheaves of objects of the gerbe
$p : \mathcal{C} \to W_{spaces, \etale}$ constructed
in Remark \ref{remark-gerbe-of-lifts} are abelian, and
\item the sheaf of groups $\mathcal{G}$ constructed in
Stacks, Lemma \ref{stacks-lemma-gerbe-abelian-auts}
is a quasi-coherent $\mathcal{O}_W$-module.
\end{enumerate}
\end{lemma}

\begin{proof}
We will prove both statements at the same time. Namely, given
an object $\xi = (U, U', a, i, x', \alpha)$ we will endow
$\mathit{Aut}(\xi)$ with the structure of a
quasi-coherent $\mathcal{O}_U$-module on $U_{spaces, \etale}$ and
we will show that this structure is compatible with pullbacks.
This will be sufficient by glueing of sheaves
(Sites, Section \ref{sites-section-glueing-sheaves})
and the construction of $\mathcal{G}$ in the proof of
Stacks, Lemma \ref{stacks-lemma-gerbe-abelian-auts}
as the glueing of the automorphism sheaves $\mathit{Aut}(\xi)$
and the fact that it suffices to check a module is
quasi-coherent after going to an \'etale covering
(Properties of Spaces, Lemma
\ref{spaces-properties-lemma-characterize-quasi-coherent}).

\medskip\noindent
We will describe the sheaf $\mathit{Aut}(\xi)$ using the
same method as used in the proof of
Lemma \ref{lemma-etale-local-lifts-isomorphic}.
Consider the algebraic space
$$
I' = U' \times_{x', \mathcal{X}', x'} U'
$$
with projections $p' : I' \to U'$ and $q' : I' \to U'$.
Over $I'$ there is a universal $2$-morphism
$\gamma : x' \circ p' \to x' \circ q'$.
The identity $x' \to x'$ defines a diagonal morphism
$$
\xymatrix{
U' \ar[rr]_{\Delta'} & & I' \ar[ld]^{p'} \ar[rd]_{q'} \\
& U' & & U'
}
$$
such that the compositions $U' \to I' \to U'$ and $U' \to I' \to U'$
are the identity morphisms. We will denote the base change of
$U', I', p', q', \Delta'$ to $\mathcal{X}$ by $U, I, p, q, \Delta$.
Since $W' \to \mathcal{X}'$ is smooth, we see that $p' : I' \to U'$
is smooth as a base change.

\medskip\noindent
A section of $\mathit{Aut}(\xi)$ over $U$ is a morphism $\delta' : U' \to I'$
such that $\delta'|_U = \Delta$ and such that
$p' \circ \delta' = \text{id}_{U'}$. To be explicit,
$(\text{id}_U, q' \circ \delta', (\delta')^*\gamma) : \xi \to \xi$
is a formula for the corresponding automorphism.
More generally, if
$f : V \to U$ is an \'etale morphism, then there is a thickening
$j : V \to V'$ and an \'etale morphism
$f' : V' \to U'$ whose restriction to $V$ is $f$ and $f^*\xi$
corresponds to $(V, V', a \circ f, j, x' \circ f', f^*\alpha)$, see proof of
Lemma \ref{lemma-gerbe-of-lifts-fibred}.
 a section of $\mathit{Aut}(\xi)$ over $V$ is a morphism
$\delta' : V' \to I'$
such that $\delta'|_V = \Delta \circ f$
and $p' \circ \delta' = f'$\footnote{A formula for the corresponding
automorphism is $(\text{id}_V, h', (\delta')^*\gamma)$.
Here $h' : V' \to V'$ is the unique (iso)morphism such that
$h'|_V = \text{id}_V$ and such that
$$
\xymatrix{
V' \ar[r]_{h'} \ar[rd]_{q' \circ \delta'} & V' \ar[d]^{f'} \\
& U'
}
$$
commutes. Uniqueness and existence of $h'$ by topological invariance of
the \'etale site, see More on Morphisms of Spaces,
Theorem \ref{spaces-more-morphisms-theorem-topological-invariance}.
The reader may feel we should instead look at morphisms
$\delta'' : V' \to V' \times_{\mathcal{X}'} V'$ with
$\delta'' \circ j = \Delta_{V'/\mathcal{X}'}$ and
$\text{pr}_1 \circ \delta'' = \text{id}_{V'}$.
This would be fine too: as $V' \times_{\mathcal{X}'} V' \to I'$ is \'etale,
the same topological invariance tells us that sending
$\delta''$ to $\delta' = (V' \times_{\mathcal{X}'} V' \to I') \circ \delta''$
is a bijection between the two
sets of morphisms.}.

\medskip\noindent
We conclude that $\mathit{Aut}(\xi)$ as a sheaf of sets agrees with
the sheaf defined in
More on Morphisms of Spaces, Remark
\ref{spaces-more-morphisms-remark-another-special-case}
for the thickenings $(U \subset U')$ and $(I \subset I')$ over
$(U \subset U')$ via $\text{id}_{U'}$ and $p'$.
The diagonal $\Delta'$ is a section of this sheaf and by
acting on this section using More on Morphisms of Spaces, Lemma
\ref{spaces-more-morphisms-lemma-action-sheaf}
we get an isomorphism
\begin{equation}
\label{equation-isomorphism}
\SheafHom_{\mathcal{O}_U}(\Delta^*\Omega_{I/U}, \mathcal{C}_{U/U'})
\longrightarrow
\mathit{Aut}(\xi)
\end{equation}
on $U_{spaces, \etale}$. There three things left to check
\begin{enumerate}
\item the construction of (\ref{equation-isomorphism})
commutes with \'etale localization,
\item $\SheafHom_{\mathcal{O}_U}(\Delta^*\Omega_{I/U}, \mathcal{C}_{U/U'})$
is a quasi-coherent module on $U$,
\item the composition in $\mathit{Aut}(\xi)$ corresponds
to addition of sections in this quasi-coherent module.
\end{enumerate}
We will check these in order.

\medskip\noindent
To see (1) we have to show that if $f : V \to U$ is \'etale,
then (\ref{equation-isomorphism}) constructed using $\xi$ over $U$,
restricts to the map (\ref{equation-isomorphism})
$$
\SheafHom_{\mathcal{O}_V}(
\Delta_V^*\Omega_{V \times_\mathcal{X} V/V}, \mathcal{C}_{V/V'}) \to
\mathit{Aut}(\xi|_V)
$$
constructed using $\xi|_V$ over $V$ on $V_{spaces, \etale}$.
This follows from the discussion in the footnote above
and More on Morphisms of Spaces, Lemma
\ref{spaces-more-morphisms-lemma-action-by-derivations-etale-localization}.

\medskip\noindent
Proof of (2). Since $p'$ is smooth, the morphism $I \to U$ is smooth,
and hence the relative module of differentials $\Omega_{I/U}$
is finite locally free (More on Morphisms of Spaces, Lemma
\ref{spaces-more-morphisms-lemma-smooth-omega-finite-locally-free}).
On the other hand, $\mathcal{C}_{U/U'}$ is quasi-coherent
(More on Morphisms of Spaces, Definition
\ref{spaces-more-morphisms-definition-conormal-sheaf}).
By Properties of Spaces, Lemma
\ref{spaces-properties-lemma-properties-quasi-coherent}
we conclude.

\medskip\noindent
Proof of (3). There exists a morphism $c' : I' \times_{p', U', q'} I' \to I'$
such that $(U', I', p', q', c')$ is a groupoid in algebraic spaces
with identity $\Delta'$. See
Algebraic Stacks, Lemma \ref{algebraic-lemma-map-space-into-stack} for example.
Composition in $\mathit{Aut}(\xi)$ is induced by the morphism
$c'$ as follows. Suppose we have two morphisms
$$
\delta'_1, \delta'_2 : U' \longrightarrow I'
$$
corresponding to sections of $\mathit{Aut}(\xi)$ over $U$ as above,
in other words, we have $\delta'_i|U = \Delta_U$ and
$p' \circ \delta'_i = \text{id}_{U'}$. Then the composition
in $\mathit{Aut}(\xi)$ is
$$
\delta'_1 \circ \delta'_2 = c'(\delta'_1 \circ q' \circ \delta'_2, \delta'_2)
$$
We omit the detailed verification\footnote{The reader can see immediately
that it is necessary to precompose
$\delta'_1$ by $q' \circ \delta'_2$ to get a well defined $U'$-valued
point of the fibre product $I' \times_{p', U', q'} I'$.}.
Thus we are in the situation described in
More on Groupoids in Spaces, Section
\ref{spaces-more-groupoids-section-groupoid-sections}
and the desired result follows from
More on Groupoids in Spaces, Lemma
\ref{spaces-more-groupoids-lemma-composition-is-addition}.
\end{proof}

\end{document}
